\subsection{Why constant-frequency restarting loses acceleration}

For a degree-$H$ polynomial method on an SPD spectrum $[\alpha,1]$, the best
worst-case block contraction is
\begin{equation}
  \rho_H^\star
  =\frac1{\abs{T_H((1+\alpha)/(1-\alpha))}},
  \label{eq:cheb-block}
\end{equation}
where $T_H$ is the Chebyshev polynomial of the first kind.  If
$H\sqrt\alpha\ll1$, then
\begin{equation}
  1-\rho_H^\star=\Theta(H^2\alpha).
  \label{eq:short-block}
\end{equation}
Repeating such blocks therefore needs
\begin{equation}
  \Omega\!\left(\frac1{H\alpha}\log\frac1\epsppr\right)
  \label{eq:restart-lower}
\end{equation}
total polynomial degree.  In particular, $H=\bigO(1)$ returns to the
nonaccelerated $\alpha^{-1}$ scale.  An accelerated accuracy epoch must retain
momentum for $H=\Omega(1/\sqrt\alpha)$ steps.  This classical polynomial fact
does not by itself resolve support changes, but it rules out constant-step
ordinary restarts as the flattening mechanism.
